% !TeX program = xelatex
% Guion de exposición para la presentación de avance (uso personal, no se entrega)
\documentclass[10pt,letterpaper]{extarticle}
\usepackage{fontspec}
\setmainfont{Inter}[BoldFont={Inter SemiBold},ItalicFont={Inter Italic},BoldItalicFont={Inter SemiBold Italic}]
\usepackage[spanish,es-noshorthands]{babel}
\usepackage[letterpaper,margin=1.7cm]{geometry}
\usepackage{parskip}
\usepackage{enumitem}
\usepackage{booktabs}
\usepackage{tabularx}
\usepackage{xcolor}
\usepackage[most]{tcolorbox}
\usepackage[hidelinks]{hyperref}
\definecolor{accent}{HTML}{0E6F7C}
\definecolor{warm}{HTML}{C2511A}
\definecolor{muted}{HTML}{66727F}
\setlist{itemsep=1pt,topsep=2pt}

\newtcolorbox{alerta}{colback=warm!7,colframe=warm,boxrule=0.6pt,arc=2pt,left=4pt,right=4pt,top=3pt,bottom=3pt}
\newtcolorbox{nota}{colback=accent!6,colframe=accent,boxrule=0.6pt,arc=2pt,left=4pt,right=4pt,top=3pt,bottom=3pt}

\newcommand{\lamina}[4]{\par\medskip{\color{accent}\large\textbf{#1 · #2}}\hfill{\small\color{muted}#3 \quad|\quad llevo #4}\par\smallskip}
\newcommand{\pasa}[1]{\par{\small\color{muted}Transición: #1}\par}

\begin{document}

{\LARGE\textbf{Guion: presentación de avance FAST-LoRa}}\\[1mm]
{\color{muted}IPD-438 · 27 de octubre de 2026 · objetivo 14 a 15 minutos (máximo permitido: 20)}

\begin{nota}
Las láminas tienen poco texto a propósito: lo que falta en ellas se dice en voz alta, y está en este guion. No hay que
leerlo; basta con tener claras las ideas y memorizar las frases de apertura y cierre. Medir el tiempo en un ensayo
completo, en voz alta y con cronómetro.
\end{nota}

\section*{Puntos de control de tiempo}
\begin{tabularx}{\linewidth}{@{}lX@{}}
\toprule
Al terminar la lámina 4 & unos 3:00. Si van más de 4:00, acortar los ejemplos de conceptos. \\
Al terminar la lámina 9 & unos 7:15. Es la mitad de la charla. \\
Al terminar la lámina 14 & unos 12:30. Si van más de 13:15, saltar la lámina del plan. \\
Cierre & entre 14:00 y 15:00. \\
\bottomrule
\end{tabularx}

\section*{Lámina por lámina}

\lamina{1}{Portada}{20 s}{0:20}
Frase memorizada para partir: “En quince minutos van a saber por qué un paquete LoRa llega o se pierde, y cómo decidir
si un simulador que calcula eso en milisegundos dice la verdad.” Sin “buenas tardes, hoy les voy a presentar”: los
primeros minutos el público presta atención gratis. El reloj es el símbolo de la charla: la vuelta completa son los
5 minutos de FLoRa y la astilla roja los 8 ms de FAST-LoRa.

\lamina{2}{Diseñar una red es elegir entre $10^{168}$ opciones}{60 s}{1:20}
Partir por el ejemplo. Cien sensores en un campo; cada uno elige un SF entre seis y una potencia entre ocho, 48 opciones
por sensor, $48^{100}$ redes posibles. Preguntar: “¿Cuánto creen que se demora evaluar una sola?” y aguantar unos
segundos de silencio. Mostrar la respuesta: más de cinco minutos en FLoRa, menos de 8 ms en FAST-LoRa. Decir la
consecuencia, que ya no está escrita: optimizar una red pide miles de evaluaciones, y diez mil son más de un mes en
FLoRa contra poco más de un minuto en FAST-LoRa.
\pasa{“Esa diferencia es lo que vende el paper. Pero hay un problema.”}

\lamina{3}{La promesa y el “sin embargo”}{70 s}{2:30}
Es la inestabilidad, en el sentido de McEnerney. A la izquierda lo que afirma el paper; explicar PDR como la fracción
de paquetes que llegan. A la derecha leer la caja naranja como una sola frase: lo validaron solo sus autores, contra un
FLoRa que ellos modificaron, con el ADR apagado (el ADR es el ajuste automático de SF y potencia, justo lo que la
herramienta promete evaluar) y con números que no calzan. Adelantar que eso se ve en el estado del arte. Cerrar con la
frase de abajo mirando al público.

\lamina{4}{Promesa y mapa}{25 s}{2:55}
Decir la promesa de forma explícita, porque la pauta la evalúa: “Al terminar van a saber cuándo creerle a un simulador
de 8 ms.” Nombrar los tres bloques y avisar que arriba a la derecha se ve en cuál vamos.

\lamina{5}{LoRaWAN en un ejemplo}{60 s}{3:55}
Ejemplo primero: el sensor de riego de la franja de abajo. Señalar los números, no la pantalla: “miren el 1”, el sensor
transmite cuando tiene algo sin escuchar el canal; es ALOHA puro, lo del curso, con un techo de 18\,\% de uso útil.
“El 2”: el mismo paquete puede llegar a dos gateways y basta con uno. “El 3”: el servidor descarta duplicados y es el
que ajusta SF y potencia, el ADR.
\pasa{“¿Y qué es el SF que tanto se repite?”}

\lamina{6}{El spreading factor}{50 s}{4:45}
Cada SF extra duplica la duración del símbolo. Con los 10 bytes que usan los autores, de SF7 a SF12 el paquete pasa
27 veces más tiempo en el aire: más energía y más choques. La sorpresa es la compensación: con SF12 el receptor entiende
una señal cien veces más débil que el ruido. Ese es el compromiso entre alcance y tiempo en el aire.

\lamina{7}{Se pierde si llega débil}{45 s}{5:30}
Primera razón para perder un paquete. La potencia que llega baja con la distancia y además varía al azar en cada
paquete; por eso no es un número sino una campana. Todo lo que queda a la derecha de la línea de sensibilidad llega:
esa área es $\psi$. Lo que queda a la izquierda se pierde.

\lamina{8}{O si choca con otro}{60 s}{6:30}
Segunda razón. Si otro paquete pisa los últimos símbolos del preámbulo y el nuestro no llega al menos 6~dB más fuerte,
se pierde. Si usan SF distintos casi siempre sobrevive, porque los SF se interfieren poco. Esa probabilidad de no chocar
es $\zeta$. La caja es el corazón del modelo: PDR es $\psi$ por $\zeta$, y con varios gateways basta que uno reciba.

\lamina{9}{Simular contra calcular}{45 s}{7:15}
La idea central, dicha tal cual: “FLoRa cuenta paquetes, FAST-LoRa calcula probabilidades.” FLoRa genera cada paquete
durante cinco días simulados, unos 432 por nodo. FAST-LoRa evalúa $\psi$ y $\zeta$ para todos los nodos de una vez. La
moneda lo resume. Marcar el cerco: FAST-LoRa no modela la bajada, varios canales ni capas superiores, y supone tráfico
estable.

\lamina{10}{Mapa del estado del arte}{60 s}{8:15}
De izquierda a derecha: la fórmula de ALOHA y los modelos de geometría estocástica, que responden por una red promedio;
los modelos por dispositivo, donde están Toro-Betancur y FAST-LoRa; los simuladores paquete a paquete, LoRaSim, ns-3 y
FLoRa; y en el extremo una red física. Más detalle cuesta más. La flecha es lo que voy a comparar. Decir el dato del pie:
FAST-LoRa tiene cuatro citas, tres del mismo grupo, y ninguna réplica independiente.

\lamina{11}{Lo que me llevo de cada lectura}{55 s}{9:10}
No leer las seis tarjetas. Contar tres: de Bor et al. salen los parámetros de propagación que usan FLoRa y FAST-LoRa,
medidos en Lancaster; Croce et al. midió que los SF se interfieren un poco, y de ahí salen los umbrales; y lo más
importante, Toro-Betancur ya había publicado este mismo modelo en 2021, y FAST-LoRa lo toma a través de Zhang et al.
Lima et al. sirve de advertencia: en otro entorno, como la Amazonía, los parámetros cambian mucho.

\lamina{12}{Cuatro cabos sueltos del paper}{70 s}{10:20}
Contarlo como historia: “Leí el paper con calculadora en mano y encontré cuatro cosas.” Uno, el tamaño del paquete no
está en el paper, y de él depende casi todo; recién en el código aparece: 10 bytes. Dos, los errores se promedian de
dos formas: el resumen da 0,040 promediando por nodo y el texto da 0,038, el promedio simple de la tabla. Tres, con sus
propios tiempos (15~s contra 0,6~ms) la aceleración es de unas diez mil veces, no mil. Cuatro, la ecuación 11 usa
$2\sigma$ donde corresponde $\sqrt2\sigma$; la demostración está en el respaldo.

\lamina{13}{Lo que muestra el código}{65 s}{11:25}
“Como faltaba el payload, fui al código.” Los experimentos no están en el repositorio que cita el paper, sino en otro
del mismo autor, con su versión modificada de FLoRa. Recorrer las cuatro filas: los dos simuladores no calculan igual el
tiempo en el aire (FLoRa usa 20 bytes fijos y suma símbolos); usan umbrales distintos para paquetes del mismo SF; en el
código la matriz de umbrales está indexada al revés; y el ADR del código falla al ejecutarse, lo que calza con que nunca
se validó con ADR. Aclarar que con el tráfico que usan algunas de estas diferencias pesan poco; medir cuánto pesa cada
una es justamente lo que aporta la réplica.

\lamina{14}{Objetivo del semestre}{50 s}{12:15}
Leer el diagrama: la misma topología entra a los dos simuladores, de ambos sale la PDR y la energía por nodo, se calcula
el error entre ellos y se compara con el del paper. El criterio de éxito quedó fijado antes de simular, para no
acomodarlo a los resultados. Si no calza, explicar por qué también es un resultado. Entregables: figuras propias sobre
las del paper, tablas de error y un repositorio público.

\lamina{15}{Resultado a reproducir}{50 s}{13:05}
“En la presentación final van a ver esta figura y esta tabla, hechas por mí.” El 1 es el caso más simple, un nodo y un
gateway (144 corridas). El 2 es la tabla de error en redes densas (80 corridas); el valor naranjo, 1,044 puntos, es el
peor caso. El 3, los tiempos de cómputo, sale de las mismas corridas.

\lamina{16}{Plan}{25 s}{13:30}
Se puede saltar. Tres fases: instalar el FLoRa de los autores y correr el experimento 1, después los experimentos 2 y 3,
y al final la comparación, el informe y el repositorio. El riesgo es que esa versión de FLoRa es antigua (1.1.0, con
OMNeT++ 6.0 e INET 4.4).

\lamina{17}{Referencias}{10 s}{13:40}
Mostrar unos segundos y decir que cada afirmación de las láminas lleva su número de referencia.

\lamina{18}{Lo que me gustaría que recuerden}{40 s}{14:20}
Repetir la idea central por tercera vez: FAST-LoRa calcula en milisegundos lo que FLoRa simula en minutos, y antes de
diseñar redes con él hay que verificarlo por fuera. Decir la frase de Hamming adaptada (creer y dudar a la vez, como
Darwin con sus notas) y nombrar los cuatro aportes. Cerrar sin “¿preguntas?” ni “gracias”, por ejemplo con “quedo atento
a sus comentarios, sobre todo si alguien ha usado FLoRa”. Si suena descortés, un “gracias” breve después no hace daño.

\section*{Si falta tiempo}
Saltar la lámina 16 y decir el plan en una frase. En la 11 contar solo la tarjeta de Toro-Betancur. En la 10 señalar
solo FAST-LoRa, FLoRa y la flecha. Nunca pasar láminas a toda velocidad.

\section*{Preguntas probables}
\begin{description}[leftmargin=0pt,itemsep=3pt]
\item[¿Por qué no Packet Tracer o GNS3?] Emulan redes IP y no modelan la capa física LoRa. Respaldo 22.
\item[¿Es LoRaWAN completo?] No: solo subida, un canal y sin ADR, igual que la validación original.
\item[¿Cómo sabes que la ecuación 11 está mal?] La resta de dos gaussianas independientes de desviación $\sigma$ tiene
varianza $2\sigma^2$, o sea desviación $\sqrt2\sigma$; el paper y el código usan $2\sigma$. Respaldo 21.
\item[¿Qué payload vas a usar?] 10 bytes, el del código de los autores. El código lo recibe en bits y FLoRa calcula el
tiempo en el aire con 20 bytes fijos sin importar lo configurado. Respaldos 20 y 23.
\item[¿Hay alguna réplica de FAST-LoRa?] No encontré ninguna: 4 citas en Semantic Scholar, 3 del mismo grupo.
\item[¿Entonces el paper está mal?] No necesariamente. Con $\lambda=0{,}001$ la probabilidad de traslape entre dos nodos
es menor a 0,2\,\%, así que varias diferencias pesan poco. Lo que falta es medir cuánto pesa cada una.
\item[¿Por qué este paper?] Es de 2026, en una revista indexada de IEEE, con código abierto: replicable en un semestre.
\end{description}

\section*{Cómo responde la charla a la pauta}
\begin{tabularx}{\linewidth}{@{}Xl@{}}
\toprule
Problema y motivación (anomalía, inestabilidad), 15 pts & 2 y 3 \\
Objetivo al término de la asignatura, 10 pts & 14 \\
Resultado que se reproducirá en la final, 5 pts & 15 \\
Contenidos previos para entender el tema, 15 pts & 5 a 9 \\
Avance en el estudio del estado del arte, 20 pts & 10 a 13 \\
Al menos dos referencias, 5 pts & citadas en cada lámina y lista en la 17 \\
Tiempo (máx. 20 min) & unos 14 a 15 minutos \\
Organización del material & mapa en la 4 e indicador de bloque arriba a la derecha \\
Uso de la tecnología & PDF 16:9; ver abajo \\
Promesa o qué recordar & promesa en la 1 y la 4; qué recordar en la 18 \\
\bottomrule
\end{tabularx}

\section*{Consejos de los tres escritos y dónde se aplican}
\begin{tabularx}{\linewidth}{@{}lX@{}}
\toprule
McEnerney & Abrir con una inestabilidad (lámina 3, “sin embargo”) y pensar en lo que el público necesita. \\
Hamming & Primero por qué importa, después los detalles (láminas 2 y 3 antes de cualquier ecuación). Creer y dudar a la
vez (frase de cierre). \\
Winston & Promesa al inicio; idea central tres veces (láminas 1, 9 y 18); números en las figuras en vez de puntero;
pregunta al público con silencio; poco texto por lámina; última lámina con lo que hay que recordar. Estrella de Winston:
símbolo (el reloj), slogan (“calcular en vez de simular”), sorpresa (decodificar bajo el ruido y los cabos sueltos),
idea saliente y la historia de cómo los encontré. \\
Peyton Jones & Primeros dos minutos en el problema; índice de tres puntos después de la motivación; ejemplo antes que
regla; pensar en el 85\,\% que no sabe de LoRa; no pedir disculpas. \\
EMBL & Cerca de una lámina por minuto; destacar el dato importante (celda naranja de la tabla); tener decidido qué
saltar. \\
\bottomrule
\end{tabularx}

\section*{Presentar por Zoom viendo las notas (sin que el público las vea)}
Se usa pdfpc con \texttt{presentacion\_avance.pdf}. Las notas están en \texttt{presentacion\_avance.pdfpc}, que pdfpc
lee solo si está en la misma carpeta. pdfpc abre dos ventanas: \emph{presentation} (solo láminas, es la que se comparte)
y \emph{presenter}, la consola del orador (lámina actual, siguiente, notas y cronómetro, que solo ves tú).
\begin{enumerate}[leftmargin=5mm]
\item Instalar una vez: \texttt{sudo apt install pdf-presenter-console}.
\item Abrir: \texttt{pdfpc -w both -d 15 presentacion\_avance.pdf} (dentro de la carpeta del proyecto). Las notas
aparecen en el panel inferior derecho de la consola; con \texttt{+} y \texttt{-} se cambia su tamaño.
\item Acomodar las ventanas: la de presentación a la izquierda, lo más grande posible, y la consola a la derecha.
\item En Zoom: Compartir pantalla, elegir la ventana de presentación de pdfpc, no la pantalla completa. Si se elige la
pantalla completa, el público ve las notas.
\item Avanzar con las flechas o la barra espaciadora, con la consola enfocada. \texttt{g} y un número salta a esa lámina
(útil para los respaldos). \texttt{b} pone la presentación en negro. \texttt{p} pausa el cronómetro.
\end{enumerate}
Tu sesión es Wayland. Si Zoom no ofrece compartir una sola ventana, cerrar sesión y entrar en la sesión Xorg (engranaje
en la pantalla de inicio de sesión), donde Zoom comparte ventanas sin problema. Probarlo antes uniéndose a la reunión con
el celular para mirar lo que ve el profesor.

Plan B, sin instalar nada: abrir \texttt{presentacion\_avance.pdf} en una ventana del visor de documentos (sin pantalla
completa) y compartir solo esa ventana, con \texttt{guion\_avance.pdf} abierto en otra ventana al lado.

\begin{alerta}
Antes de presentar, revisar: (1) la fecha y el código del curso de la portada; (2) que el plan de la lámina 16 refleje
lo que realmente llevas instalado; (3) que puedas explicar con tus palabras las láminas 12 y 13, revisando el paper,
los respaldos 20 y 21 y los repositorios (github.com/jaznrp/FastLoRa y github.com/jaznrp/MultiAgent-MARLTC).
\end{alerta}

\end{document}
